\documentclass[12pt]{article}
\usepackage[utf8]{inputenc}
\usepackage[T1]{fontenc}
\usepackage{hyperref}
\title{Yopno Language Resources for K-12 Curriculum Development}
\author{Compiled by Claude Sonnet 5.5 (claude-sonnet-5-5) for Andr\'as Kornai}
\date{}
\begin{document}
\maketitle

\section{Introduction}
Yopno (ISO 639-3 yut, Glottocode yopn1238) is a Trans-New Guinea language of Papua New Guinea with about 9,000 speakers. Most live along the Yopno river and its tributaries in the Finisterre Mountains, in Morobe and Madang provinces. This report rests on a single search, and the only resources found were Scripture texts.

\section{Primary Educational Resources}
No primers, textbooks or school materials were found.

\section{Linguistic and Reference Materials}
No dictionary or grammar was found in the results.

\section{Religious and Community Resources}
The New Testament and Psalms, \emph{Ping Kop Gen}, were published in 2010 by Wycliffe Bible Translators, with a web version in 2015. The licence allows sharing but not sale or alteration of the text.
\begin{itemize}
\item eBible web edition: \url{https://ebible.org/yut/FRT01.htm}
\item Full PDF: \url{https://ebible.org/pdf/yut/yut_all.pdf}
\item Scripture Earth PDF: \url{https://www.scriptureearth.org/data/yut/PDF/00-PByut-web.pdf}
\item EPUB: \url{https://www.ebible.org/epub/yut.epub}
\end{itemize}

\section{Digital and Audio Resources}
None found beyond the electronic Scripture editions above.

\section{Curriculum Development Context}
Yopno is spoken by all members of its community, which is a strong basis for mother-tongue education. The only documented text is the New Testament, which could serve as a source of written text and for orthography study. Everything else would need to be created or located through further searching.
\end{document}